\label{sec:algo_select}

All six algorithms described in Section~\ref{sec:alg_overview} were
evaluated. This section explains the rationale behind each inclusion and
details the primary algorithm, AOS.

\textit{GPS as a deliberate baseline.} The central empirical question is
whether population-based global search is necessary on a moving-average (MA)
parameter fitness landscape, or whether a single-state local optimiser is
sufficient. GPS~\cite{torczon1997} has no population, no escape mechanism
from local optima, and no memory of previously evaluated positions beyond the
current best. Including it provides a direct lower bound: if GPS performs
competitively, population diversity contributes little; if GPS underperforms
consistently, the landscape demands it.

\textit{PSO as a population-based benchmark.} PSO~\cite{shi1998} is one of
the most extensively studied metaheuristics and provides a reliable
population-based performance reference. It is also structurally related to
AOS --- both use a global nucleus/best to guide the population --- making a
paired comparison directly informative about the contribution of personal
best memory. If AOS and PSO are statistically indistinguishable, per-particle
historical memory may or may not add marginal benefit in this problem domain.

\textit{AOS as the physics based algorithm compared to the others being 
biology based.} The MA parameter fitness landscape is
expected to be multimodal: distinct parameter combinations can produce
similar short-term profits through qualitatively different mechanisms,
creating multiple local peaks of comparable height. AOS~\cite{azizi2021} was
selected because its shell-based population structure explicitly supports
simultaneous exploitation of multiple landscape regions. Candidate solutions
(electrons) are assigned to concentric shells around the nucleus (current
global best) according to their fitness rank; electrons in inner shells
exploit promising regions while outer-shell electrons conduct broader
exploration. Transitions between shells are governed by a quantum-mechanical
photon absorption and emission analogy: an electron absorbs energy (moves to
an outer shell, exploration) when its transition energy falls below the
shell's binding energy, and emits energy (moves to an inner shell,
exploitation) otherwise. This structured push-pull between shells provides
more persistent population spread than PSO's convergent swarm behaviour,
which collapses toward the global best over time. AOS was published in 2021
and has not, to our knowledge, been applied to financial parameter
optimisation in the literature.

The pseudocode below summarises the core AOS iteration:

%TC:ignore
\begin{lstlisting}[basicstyle=\small\ttfamily, breaklines=true, columns=fullflexible, keepspaces=true, frame=single, xleftmargin=1.5em, xrightmargin=1em]
Initialise N electrons randomly in the search space
Set nucleus := position of the best electron (highest fitness)
Assign electrons to K shells by fitness rank (binding energy)

FOR each iteration t = 1, 2, ..., T:
    FOR each shell k = 1, ..., K:
        Compute shell binding energy E_k
        FOR each electron i in shell k:
            Sample transition energy E
            IF E < E_k  (electron absorbs energy):
                Move to outer shell; apply exploration step
            ELSE  (electron emits energy):
                Move to inner shell; apply exploitation step
            Update position via quantum transition equation
            Evaluate fitness of new position
    Apply LE perturbation: randomly displace the
        lowest-energy electron to escape local optima
    Update nucleus if any electron improves best fitness

RETURN nucleus position as best solution found
\end{lstlisting}
%TC:endignore

\textit{LE degradation --- an observed weakness.} The Lowest Energy (LE)
perturbation step is designed to prevent stagnation by randomly displacing
the electron with the worst fitness. In practice, the LE designation is not
stable: as other electrons improve their fitness across iterations, the
electron currently labelled LE may no longer occupy the position most in need
of escape relative to the updated nucleus. The original
paper~\cite{azizi2021} does not address this drift. This was observed
empirically across runs and is logged per iteration; it is documented here
rather than patched, as any correction would constitute a modification to the
published algorithm.

\textit{Metaparameter choices.} Population size was scaled with
dimensionality --- larger populations were used for the $d=14$ weighted
strategy to partially offset the curse of dimensionality. The shell count
$n_\text{shells}$ is implemented as an upper cap rather than a fixed
assignment: each iteration uses $\min(n_\text{shells}, |\text{remaining
population}|)$ shells to prevent degenerate empty-shell states that can arise
in high-dimensional runs where effective population diversity is reduced. The
iteration budget was fixed uniformly across all six algorithms to ensure a
controlled comparison of per-iteration effectiveness.
